\subsection{中山引理}

命题 6. —— 对每个 A-模 M，我们有 \(\Re(A)M \subset \Re(M)\)。若 A-模 M 是投射的，则有等式。

设 P 是 M 的一个极大子模；A-模 M/P 是单的，故由 VIII, 第 155 页的命题 5，它被 \(\Re(\mathrm{A})\) 零化。于是对 M 的每个极大子模 P，我们有 \(\Re(\mathrm{A})\mathrm{M}\subset\mathrm{P}\)，因而 \(\Re(\mathrm{A})\mathrm{M}\subset\Re(\mathrm{M})\)。

我们显然有 \(\mathfrak{R}(A_{s})=\mathfrak{R}(A)A_{s}\)。若 A-模 M 是投射的，则存在一个 A-模 N，使得 \(M\oplus N\) 是自由的，即是模的一个族 \((\mathrm{L}_{i})_{i\in\mathrm{I}}\)（同构于 \(A_{s}\)）的直和。由 VIII, 第 152 页的推论 2，我们有 \(\mathfrak{R}(M\oplus N)=\mathfrak{R}(M)\oplus\mathfrak{R}(N)\) 与 \(\mathfrak{R}\big(\bigoplus_{i\in I}L_{i}\big)=\bigoplus_{i\in I}\mathfrak{R}(L_{i})\)。于是，\(\mathfrak{R}(M)=\mathfrak{R}(A)M\) 由等式 \(\mathfrak{R}(L_{i})=\mathfrak{R}(A)L_{i}\) 推出。

定理 2（``中山引理''）. —— 设 M 是一个 A-模，a 是 A 的一个双边理想。设下列条件之一满足：

\begin{enumerate}
\def\labelenumi{(\roman{enumi})}
\item
  A-模 M 是有限生成的，且 \(\mathfrak{a}\) 包含于 A 的根中。
\item
  理想 \(\mathfrak{a}\) 是幂零的。
\end{enumerate}

若 N 是 M 的一个子模，使得 \(M = N + aM\)，则我们有 N = M。特别地，若模 M 非零，则我们有 \(M \neq aM\)。

设 M 是有限生成的，且我们有 \(\mathfrak{a} \subset \Re(\mathrm{A})\)。设 N 是 M 的一个子模，使得 \(M = N + aM\)。由命题 6，我们有 \(M = N + \Re(M)\)，故由 VIII, 第 153 页的命题 2, b) 得 M = N。

现在设 a 是幂零的，并设 N 是 M 的一个子模，使得 \(M = N + aM\)。对整数 \(n \geqslant 0\) 作归纳，我们建立关系 \(M = N + a^{n}M\)。按假设，存在一个整数 \(n \geqslant 0\) 使得 \(a^{n} = 0\)；因而 M = N。

定理的最后一个断言由上述取 N 为 0 即得。

推论 1. —— 我们保持定理 2 的假设。设 \((x_{i})_{i\in\mathrm{I}}\) 是 M 的元素的一个族，\(\overline{x}_{i}\) 是 \(x_{i}\) 在 M/aM 中的典范像。若族 \((\overline{x}_{i})_{i\in\mathrm{I}}\) 生成 (A/a)-模 M/aM，则族 \((x_{i})_{i\in\mathrm{I}}\) 生成 A-模 M。

这由把定理 2 应用于由族 \((x_{i})_{i\in \mathbf{I}}\) 生成的 M 的子模 N 得出。

推论 2. —— 我们保持定理 2 的假设。此外，设 M' 是一个 A-模，u : M' → M 是一个同态。若由 u 通过取商得到的从 M'/aM' 到 M/aM 的同态 \(\overline{u}\) 是满射，则 u 是满射。

只需把定理 2 应用于 u 的像 N：事实上，\(\overline{u}\) 的像是 \((\mathrm{N} + \mathfrak{aM})/\mathfrak{aM}\)，故 \(\overline{u}\) 是满射当且仅当我们有 \(N + aM = M\)。

